\SourceLeaf{8}
\pdfbookmark[0]{Investigation of the Elliptic Elements of Pallas}{palladis}
\begin{center}
{\large INVESTIGATION}\par\medskip
{\Large OF THE ELLIPTIC ELEMENTS OF PALLAS}\par\medskip
FROM THE OPPOSITIONS OF 1803, 1804, 1805, 1807, 1808, 1809.
\end{center}
\Article{1}
The primary planet Pallas, discovered in March 1802, has had seven periodic apparitions so far. But since the planet had already passed opposition with the Sun on the day it was first discovered, only six oppositions have been counted to date; most of these were observed fairly well at several observatories and could therefore be determined with sufficient precision. The opposition of 1808 alone must be excepted. In that year the planet was at the aphelion of its orbit, and its faint light caused considerable difficulties for observers. To tell the truth, astronomers neglected it somewhat. The distinguished \textsc{von Lindenau} did provide quite a large number of right ascensions, observed with his usual care around the time of opposition with the excellent meridian telescope of the Seeberg Observatory. From elsewhere, however, I could obtain only a few declinations, and these were much less certain. It was therefore impossible to establish this opposition with the same precision as the others, and the latitude in particular remained too uncertain.
\Article{2}
As soon as the observations for any apparition had ended, I made a practice each year of collecting them all, examining them carefully, and combining them with the observations of preceding years so as to obtain a still \SourceLeaf{9}more accurate determination of the elliptic elements. The last calculation of this kind, carried out toward the end of 1807, yielded the elements published in the January 1808 issue of Vol. XVII of the distinguished \textsc{von Zach}’s Commercium literarium astronomicum. They agreed very well with all four oppositions observed up to that time. The next year, after the fifth opposition, the scarcity of reliable observations deterred me from making new calculations. I returned to them only when the observations made in 1809 around the sixth opposition allowed me to undertake the calculation of both oppositions together. I announced it in our Learned Notices of 24 February 1810.

Here is an overview of all six oppositions observed so far:
\begin{center}\small\setlength{\tabcolsep}{3pt}
\begin{tabular}{l r r|r|r|r}
\multicolumn{3}{c|}{\shortstack{Time of opposition at the\\Göttingen meridian}}&\shortstack{Days from the start\\of 1803}&\shortstack{Heliocentric\\longitude}&\shortstack{Geocentric\\latitude}\\\hline
1803 June&30.&$0^{\mathrm h}\;27^{\mathrm m}\;32^{\mathrm s}$&181.019120&$277^0\;39'\;24''0$&$+46^0\;26'\;36''0$\\
1804 Aug.&30.&$4\;58\;27$&608.207257&$337\;0\;36.1$&$+15\;1\;49.8$\\
1805 Nov.&29.&$11\;15\;4$&1064.468796&$67\;20\;42.9$&$-54\;30\;54.9$\\
1807 May&4.&$14\;37\;41$&1585.609502&$223\;37\;27.7$&$+42\;11\;25.6$\\
1808 July&26.&$21\;17\;32$&2034.887176&$304\;2\;59.7$&$+37\;43\;53.7$\\
1809 Sept.&22.&$16\;10\;20$&2457.673843&$359\;40\;4.4$&$-7\;22\;10.1$
\end{tabular}\end{center}
\Article{3}
If Pallas moved exactly in an ellipse and followed \textsc{Kepler}’s laws, the latest elements, which agree very well with the four oppositions of 1803, 1804, 1805, 1807, could still contain only very slight errors, and they would surely have to agree with the next oppositions as well to within one or two minutes of arc. In fact the agreement fell so far short of this that those elements already showed a discrepancy of up to four minutes of arc at the fifth opposition and were twelve minutes of arc off at the sixth. Our planet is subject to such large perturbations from the other planets, especially Jupiter, that exact and lasting agreement between the phenomena and purely elliptic motion cannot be obtained. It is also clear from this that different elements will result according to which four oppositions are taken as the basis. This was confirmed by calculating two other sets of elements, one of which I recently derived from the oppositions \SourceLeaf{10}of 1804, 1805, 1807, 1808, and the other from those of 1805, 1807, 1808, 1809. To make it easier to compare these different sets of elements, I give each of them here.

\Needspace{14\baselineskip}
\begin{center}I.\quad\textit{Elliptic elements of Pallas from the oppositions of}\par
1803, 1804, 1805, 1807.\end{center}
\begingroup\small\setlength{\tabcolsep}{2pt}
\noindent\begin{tabular}{@{}p{.51\linewidth}lr@{}}
Epoch of mean longitude for 1803 at the Göttingen meridian&. . . . . . . . . . . . . .&$221^0\;39'\;30''4$\\
Mean daily tropical motion&. . . . . . . . . . . .&$770''2143$\\
Longitude of perihelion 1803&. . . . . . . . . . . . .&$121^0\;3'\;11''4$\\
Longitude of ascending node 1803&. . . . . . . . . . .&$172\;28\;56.9$\\
Inclination of orbit&. . . . . . . . . . . . . . . .&$34\;37\;41.0$\\
Eccentricity $(=\sin.\,14^0\;10'\;58''81)$&. . . . . . . . . .&$0.2450198$\\
Logarithm of semimajor axis&. . . . . . . . . . . . .&$0.4423149$
\end{tabular}\par\endgroup

\Needspace{13\baselineskip}
\begin{center}II.\quad\textit{Elliptic elements of Pallas from the oppositions of}\par
1804, 1805, 1807, 1808.\end{center}
\begingroup\small\setlength{\tabcolsep}{2pt}
\noindent\begin{tabular}{@{}p{.51\linewidth}lr@{}}
Epoch of mean longitude 1803&. . . . . . . . . . .&$221^0\;34'\;56''7$\\
Mean tropical daily motion&. . . . . . . . . . . .&$770''4467$\\
Longitude of perihelion 1803&. . . . . . . . . . . . .&$121^0\;5'\;22''1$\\
Longitude of ascending node 1803&. . . . . . . . . . .&$172\;28\;46.8$\\
Inclination of orbit&. . . . . . . . . . . . . . . .&$34\;37\;31.5$\\
Eccentricity $(=\sin.\,14^0\;10'\;4''08)$&. . . . . . . . . . .&$0.2447624$\\
Logarithm of semimajor axis&. . . . . . . . . . . . .&$0.4422276$
\end{tabular}\par\endgroup

\Needspace{13\baselineskip}
\begin{center}III.\quad\textit{Elliptic elements of Pallas from the oppositions of}\par
1805, 1807, 1808, 1809.\end{center}
\begingroup\small\setlength{\tabcolsep}{2pt}
\noindent\begin{tabular}{@{}p{.51\linewidth}lr@{}}
Epoch of mean longitude 1803&. . . . . . . . . . .&$221^0\;23'\;24''6$\\
Mean tropical daily motion&. . . . . . . . . . . .&$770''9265$\\
Longitude of perihelion 1803&. . . . . . . . . . . . .&$120^0\;58'\;4''8$\\
Longitude of ascending node 1803&. . . . . . . . . . .&$172\;27\;52.4$\\
Inclination of orbit&. . . . . . . . . . . . . . . .&$34\;36\;49.4$\\
Eccentricity $(=\sin.\,14^0\;9'\;36''63)$&. . . . . . . . . . .&$0.2446335$\\
Logarithm of semimajor axis&. . . . . . . . . . . . .&$0.4420473$
\end{tabular}\par\endgroup
\SourceLeaf{11}
\Article{4}
I showed in Section II of Book II of the Theoria motus corporum coelestium how the elliptic orbit of a planet can be determined from four observations, only two of which are complete. The solution given there treats the problem in the greatest generality. For the special case in which the four observations are oppositions, however, it is preferable to use another method. I hope astronomers will welcome an account of it. This method does require approximate knowledge of each of the elements, or at least of the four given by the inclination of the orbit, the longitude of the node, the longitude of the perihelion, and the eccentricity. But nothing prevents us from assuming that we already have this knowledge from an earlier calculation. The essential point depends on solving the problem below.
\begin{quote}\itshape
‘Given four longitudes of the planet in its orbit corresponding to given times, find the longitude of the perihelion, the eccentricity, the mean daily motion, and the epoch of mean longitude’
\end{quote}
I shall give the solution in the next article.
\Article{5}
Since the longitude of the perihelion is assumed to be approximately known already, subtract it from the longitudes in the orbit to obtain the approximate true anomalies $v$, $v'$, $v''$, $v'''$. That longitude may either be treated as constant or, if its annual variation is known, adjusted to each time as appropriate. Let $e=\sin\varphi$ be the approximate eccentricity. Using it, compute by the known methods the mean anomalies $M$, $M'$, $M''$, $M'''$ corresponding to the true anomalies $v$, $v'$, $v''$, $v'''$. Now denote the times by $t$, $t'$, $t''$, $t'''$. Their differences $t'-t$, $t''-t'$, $t'''-t''$ would have to be proportional to the differences $M'-M$, $M''-M'$, $M'''-M''$ if we had started from the true values of the longitude of the perihelion and of the angle $\varphi$. Otherwise, the corrections to be applied to these quantities will be found as follows. On setting
\[
\begin{aligned}
m&=-\frac{\cos\varphi^3}{(1+e\cos v)^2},&n&=\frac{m(2+e\cos v)\sin v}{\cos\varphi'}\\[6pt]
m'&=-\frac{\cos\varphi^3}{(1+e\cos v')^2},&n'&=\frac{m'(2+e\cos v')\sin v'}{\cos\varphi}\\[6pt]
m''&=-\frac{\cos\varphi^3}{(1+e\cos v'')^2},&n''&=\frac{m''(2+e\cos v'')\sin v''}{\cos\varphi}\\[6pt]
m'''&=-\frac{\cos\varphi^3}{(1+e\cos v''')^2},&n'''&=\frac{m'''(2+e\cos v''')\sin v'''}{\cos\varphi}
\end{aligned}
\]
\noindent\SourceLeaf{12}
it follows from Article 15 of the Theoria Motus Corporum Coelestium that if we imagine the perihelion corrected by the small increment $\mathrm d\Pi$ and the angle $\varphi$ by the increment $\mathrm d\varphi$, the corrected values of the mean anomalies become
\[
\begin{gathered}
M+m\mathrm d\Pi+n\mathrm d\varphi\\
M'+m'\mathrm d\Pi+n'\mathrm d\varphi\\
M''+m''\mathrm d\Pi+n''\mathrm d\varphi\\
M'''+m'''\mathrm d\Pi+n'''\mathrm d\varphi
\end{gathered}
\]
provided that powers and products of the corrections $\mathrm d\Pi$, $\mathrm d\varphi$ can be neglected. We thus obtain three expressions for the mean daily motion relative to the perihelion, namely
\[
\begin{gathered}
\frac{M'-M}{t'-t}+\frac{m'-m}{t'-t}\mathrm d\Pi+\frac{n'-n}{t'-t}\mathrm d\varphi\\[5pt]
\frac{M''-M'}{t''-t'}+\frac{m''-m'}{t''-t'}\mathrm d\Pi+\frac{n''-n'}{t''-t'}\mathrm d\varphi\\[5pt]
\frac{M'''-M''}{t'''-t''}+\frac{m'''-m''}{t'''-t''}\mathrm d\Pi+\frac{n'''-n''}{t'''-t''}\mathrm d\varphi
\end{gathered}
\]
from whose equality the unknowns $\mathrm d\Pi$, $\mathrm d\varphi$ will be obtained. Substituting them into one of the expressions yields the corrected mean motion, from which the semimajor axis too will be determined by the known method. Finally, substituting the values of the corrections $\mathrm d\Pi$, $\mathrm d\varphi$ into the expression for any one of the mean anomalies gives its corrected value and hence the epoch of mean longitude, to be transferred to an arbitrary time.

It is also easy to see how account may be taken of both periodic and secular perturbations, if necessary. The given longitudes in the orbit need only be cleared of the periodic perturbations before they enter the calculation. The secular perturbations must be used to transfer the longitude of the perihelion and the eccentricity assumed for an arbitrary epoch to the individual times $t$, $t'$, $t''$, $t'''$.
\Article{6}
Let us now see how this can be used to determine the orbit from four observed oppositions. At the first opposition, let the longitude of the planet in the ecliptic be $\alpha$, the geocentric latitude $\beta$, the distance of the Earth from the Sun $R$, the approximate longitude of the ascending node $\AscendingNode$, and the approximate inclination of the orbital plane to the ecliptic $i$. I denote the quantities for the other oppositions by similar symbols distinguished by one, two, or three indices. From $\alpha$, $\AscendingNode$, $i$ we obtain the longitude in the orbit by the known
\SourceLeaf{13}formula, and a similar calculation is based on the other three longitudes in the ecliptic. From these four longitudes in the orbit we obtain the elliptic elements by the instructions in the preceding article, and from these in turn the four radius vectors $r$, $r'$, $r''$, $r'''$ by the known formulas. Now set
\[
\frac{R\sin\beta}{r}=\sin(\beta-\gamma),\qquad \frac{R'\sin\beta'}{r'}=\sin(\beta'-\gamma')\quad\text{etc.}
\]
Then $\gamma$, $\gamma'$, $\gamma''$, $\gamma'''$ will evidently be the four heliocentric longitudes derived from the geocentric latitudes. Let us derive them also from the longitudes by setting
\[
\operatorname{tang}i\sin(\alpha-\AscendingNode)=\operatorname{tang}\delta,\qquad \operatorname{tang}i\sin(\alpha'-\AscendingNode)=\operatorname{tang}\delta'\ \text{etc.}
\]
It is clear that $\gamma=\delta$, $\gamma'=\delta'$, $\gamma''=\delta''$, $\gamma'''=\delta'''$ would have to hold if the true values had been taken for $\AscendingNode$ and $i$, provided that the same ellipse could satisfy all four oppositions, which make the problem overdetermined. If it turns out otherwise, the corrections $\mathrm d\AscendingNode$, $\mathrm di$ to be applied to $\AscendingNode$, $i$ are obtained as follows. Insofar as these corrections are treated as small quantities of the first order, whose powers and products may be neglected, the values $\AscendingNode+\mathrm d\AscendingNode$, $i+\mathrm di$ will evidently yield, in place of the latitudes $\gamma$, $\delta$, $\gamma'$, $\delta'$ and so on, others of this form
\[
\begin{array}{ll}
\gamma+a\mathrm d\AscendingNode+b\mathrm di,&\delta+c\mathrm d\AscendingNode+f\mathrm di\\
\gamma'+a'\mathrm d\AscendingNode+b'\mathrm di,&\delta'+c'\mathrm d\AscendingNode+f'\mathrm di\quad\text{etc.}
\end{array}
\]
where the coefficients $c$, $f$, $c'$, $f'$ and so on are found from the formulas
\[
\begin{array}{ll}
c=-\frac12\sin2\delta\operatorname{cotang}(\alpha-\AscendingNode),&f=\dfrac{\sin2\delta}{\sin2i}\\[5pt]
c'=-\frac12\sin2\delta'\operatorname{cotang}(\alpha'-\AscendingNode),&f'=\dfrac{\sin2\delta'}{\sin2i}\quad\text{etc.}
\end{array}
\]
We could also determine the coefficients $a$, $b$ and so on by analytical operations, but I prefer the following method. I repeat the calculation just described, replacing the longitude of the node $\AscendingNode$ by another value that differs from it by a small amount chosen at will. This readily gives the numerical values of the coefficients $a$, $a'$ and so on. Repeating the same calculation again with a slightly different inclination similarly gives the values of the coefficients $b$, $b'$ and so on. At the same time, this shows how large the corrections to the elliptic elements will be when the quantities $\AscendingNode$, $i$ are corrected by small amounts.

\SourceLeaf{14}
This gives \textit{four} equations for determining the corrections $\mathrm d\AscendingNode$, $\mathrm di$
\[
\begin{gathered}
\gamma-\delta+(a-c)\mathrm d\AscendingNode+(b-f)\mathrm di=0\\
\gamma'-\delta'+(a'-c')\mathrm d\AscendingNode+(b'-f')\mathrm di=0\\
\gamma''-\delta''+(a''-c'')\mathrm d\AscendingNode+(b''-f'')\mathrm di=0\\
\gamma'''-\delta'''+(a'''-c''')\mathrm d\AscendingNode+(b'''-f''')\mathrm di=0
\end{gathered}
\]
The unavoidable errors in the observations will hardly allow these equations to be satisfied exactly. The most suitable values of the unknowns $\mathrm d\AscendingNode$, $\mathrm di$ will therefore be determined by the principles explained in Section III of Book II of the Theoria Motus Corporum Coelestium. From these values the corrections to the other elements are readily derived as just described,
\Article{7}
Let me add a few observations to these instructions.

I. It hardly needs pointing out that our formulas also apply when the nodes are not regarded as stationary. One need only reduce the assumed longitude of the node from its epoch to the four times as appropriate and use these four different values for the respective four positions in those formulas. The same holds for the inclination, provided that its secular variation is known and seems worth taking into account. If we also wish to include the perturbations of the heliocentric latitude in the calculation, they must evidently be either subtracted from $\gamma$, $\gamma'$, $\gamma''$, $\gamma'''$ or added to $\delta$, $\delta'$, $\delta''$, $\delta'''$ in the preceding equations.

III. Whenever the orbital inclination and eccentricity are moderate, the coefficients $a$, $b$, $a'$, $b'$ and so on become so small that they can be neglected and the second and third calculations prescribed above can therefore be dispensed with. After finding the corrections $\mathrm d\AscendingNode$, $\mathrm di$, unless they happen to be very small, one must then calculate the longitudes in the orbit afresh from the corrected values of the longitude of the node and the inclination, namely $\AscendingNode+\mathrm d\AscendingNode$, $i+\mathrm di$. If the corrected longitudes in the orbit differ from the earlier uncorrected values by several seconds of arc, the elliptic elements must be calculated again according to Article 5. In any case it will hardly ever be necessary to recalculate the coefficients $m$, $n$, $m'$ and so on, as their values are already found accurately enough under the first hypothesis.
\SourceLeaf{15}
\Article{8}
To make the preceding instructions clearer, I give here the complete calculation by which the third set of elements was determined from the oppositions of 1805, 1807, 1808, 1809. I assumed the longitude of the ascending node for the beginning of 1803 to be $172^0\,28'\,46''8=\AscendingNode$ and the orbital inclination to be $34^0\,37'\,31''5$. I reduced the longitude to the times of the individual oppositions by adding the precession $2'\,26''01$, $3'\,37''50$, $4'\,39''12$, $5'\,37''11$. I then used the longitudes to derive the heliocentric latitudes
\[
\begin{array}{rrr}
\delta=-33^0&40'&50''63\\
\delta'=+28&14&51.24\\
\delta''=+27&20&55.86\\
\delta'''=-4&52&28.44
\end{array}
\]
\[
\begin{array}{ll}
c=+0.1252,&f=-0.9870\\
c'=-0.3366.&f'=+0.8917\\
c''=+0.3609.&f''=+0.8727\\
c'''=+0.6803,&f'''=-0{,}1811
\end{array}
\]
We also obtain the arguments of latitude
\[
\begin{array}{rrr}
257^0&25'&6''73\\
56&24&5.96\\
126&2&55.07\\
188&36&2.69
\end{array}
\]
from which the elliptic elements are to be obtained. For this purpose I set the angle $\varphi=14^0\,10'\,4''08$, as it was found in the second set of elements, and the longitude of the perihelion for the beginning of 1803 at $=121^0\,5'\,22''1$. The perihelion, like the node, is assumed to be fixed relative to the fixed stars, so its distance from the ascending node is constant, $=308^0\,36'\,35''3$. The known formulas then give the following.
\[
\begin{array}{rrr@{\quad}rrr}
v=308^0&48'&31''43,&M=328^0&15'&45''08\\
v'=107&47&30.66,&M'=79&46&27.05\\
v''=177&26&19.77,&M''=175&54&28.87\\
v'''=239&59&27.39,&M'''=266&29&57.59
\end{array}
\]
\SourceLeaf{16}
\[
\begin{array}{ll}
m=-0.68517,&n=+1.18569\\
m'=-1.06482,&n=-2.01318\\
m''=-1.59701,&n=-0.12921\\
m'''=-1.18352,&n=+1.93947
\end{array}
\]
Denoting the correction to the longitude of the perihelion by $\mathrm d\Pi$ and the correction to the angle $\varphi$ by $\mathrm d\varphi$, we obtain the mean sidereal motion

from the first to the second opposition, in $521.140706$ days
\[
111^0\,30'\,41''97-0.37965\mathrm d\Pi-3.19887\mathrm d\varphi
\]
from the second to the third opposition, in $449.277674$ days
\[
96^0\,8'\,1''82-0.53219\mathrm d\Pi+1.88397\mathrm d\varphi
\]
from the third to the fourth opposition, in $422.786667$ days
\[
90^0\,35'\,28''72+0.41349\mathrm d\Pi+2.06868\mathrm d\varphi
\]
This gives three expressions for the mean daily sidereal motion
\[
\begin{gathered}
770''31398-0.0007285\mathrm d\Pi-0.0061382\mathrm d\varphi\\
770.30718-0.0011845\mathrm d\Pi+0.0041933\mathrm d\varphi\\
771.37892+0.0009780\mathrm d\Pi+0.0048930\mathrm d\varphi
\end{gathered}
\]
whose comparison gives the following equations
\[
\begin{gathered}
0=0''00680+0.0004560\mathrm d\Pi-0.0103315\mathrm d\varphi\\
0=-1.07174-0.0021625\mathrm d\Pi-0.0006997\mathrm d\varphi
\end{gathered}
\]
from which $\mathrm d\Pi=-488''82$, $\mathrm d\varphi=-20''92$, and hence
\begin{center}\fontsize{9}{11}\selectfont\setlength{\tabcolsep}{3pt}\begin{tabular}{lr}
corrected longitude of the perihelion at the beginning of 1803,&$120^0\,57'\,13''28$,\\
and corrected value of the angle $\ .\ .\ .\ .\ .\ .$&$\varphi=14^0\,9'\,43''16$\\
the mean daily sidereal motion then becomes $\ .\ .\ .\ .\ .$&$770''7985$\\
corrected mean anomaly for the first opposition $\ .$&$328^0\,20'\,55''20$\\
mean longitude for the same epoch $\ .\ .\ .\ .\ .$&$89^0\,20'\,34''49$
\end{tabular}\end{center}
Finally the mean sidereal motion gives the logarithm of the semimajor axis $0.4420439$, and hence the logarithm of the semiparameter $0.4152361$. Now that we have the corrected values of the true anomalies
\SourceLeaf{17}
\[
\begin{array}{rrr}
v=308^0&56'&40''25\\
v'=107&55&39.48\\
v''=177&34&28.59\\
v'''=240&7&36.21
\end{array}
\]
we calculate the logarithms of the radius vectors
\[
\begin{aligned}
\log r&=0.3531088\\
\log r'&=0.4492406\\
\log r''&=0.5369700\\
\log r'''&=0.4716739
\end{aligned}
\]
From the solar tables we also have
\[
\begin{aligned}
\log R&=9.9937332\\
\log R'&=0.0039862\\
\log R''&=0.0065917\\
\log R'''&=0.0011160
\end{aligned}
\]
from which we finally deduce
\[
\begin{array}{rrr}
\gamma=-33^0&39'&48''15\\
\gamma'=+28&15&0.73\\
\gamma''=+27&20&9.07\\
\gamma'''=-4&52&53.99
\end{array}
\]

To obtain the coefficients $a$, $b$ and so on, I form a \textit{second hypothesis}, retaining the inclination but increasing the longitude of the node by one minute of arc, so that at the beginning of 1803 it is $172^0\,29'\,46''8$. Then, setting the longitude of the perihelion for the same epoch at $=121^0\,5'\,22''1$ as in the first hypothesis, we find the true anomalies
\[
\begin{array}{rrr}
v=308^0&48'&40''93\\
v'=107&47&34.06\\
v''=177&26&22.25\\
v'''=239&59&15.00
\end{array}
\]
and from these, also setting $\varphi=14^0\,10'\,4''08$, the mean anomalies
\SourceLeaf{18}
\[
\begin{array}{rrr}
M=328^0&15'&51''59\\
M'=79&46&30.67\\
M''=175&54&32.83\\
M'''=266&29&42.93
\end{array}
\]
(For this calculation the usual method is unnecessary. It suffices to add to the values of the mean anomalies found under the first hypothesis the products of the coefficients $m$, $m'$, $m''$, $m'''$, taken as positive, and the differences between the corresponding values of the true anomalies in the two hypotheses, namely
\[
\begin{aligned}
+9''50\times0.68517&=+6''51\\
+3.40\times1.06482&=+3.62\\
+2.48\times1.59701&=+3.96\\
-12.39\times1.18352&=-14.66)
\end{aligned}
\]
\Needspace{6\baselineskip}
Keeping the values of the coefficients $m$, $n$ and so on then gives
\begin{center}\fontsize{9}{11}\selectfont\setlength{\tabcolsep}{3pt}\begin{tabular}{p{.23\linewidth}p{.70\linewidth}}
$\mathrm d\Pi=-468''21$,&longitude of the perihelion at the beginning of 1803, $120^0\,57'\,33''89$\\
$\mathrm d\varphi=-20''62$,&$\varphi=14^0\,9'\,43''46$
\end{tabular}\end{center}
\begin{center}\fontsize{9}{11}\selectfont\setlength{\tabcolsep}{3pt}\begin{tabular}{ll}
mean daily sidereal motion&$=770''7761$\\
logarithm of the semimajor axis&$=0.4420523$\\
mean longitude at the first opposition&$=89^0\,20'\,47''85$
\end{tabular}\end{center}
\[
\begin{array}{rrr}
\gamma=-33^0&39'&51''10\\
\gamma'=+28&15&1.27\\
\gamma''=+27&20&10.97\\
\gamma'''=-4&52&54.66
\end{array}
\]

Finally I form a \textit{third hypothesis} by setting $\AscendingNode=172^0\,28'\,46''6$, $i=34^0\,38'\,31''5$. From this we deduce, just as before,
\begin{center}\fontsize{9}{11}\selectfont\setlength{\tabcolsep}{3pt}\begin{tabular}{lr}
longitude of the perihelion for the beginning of 1803,&$120^0\,55'\,34''46$\\
&$\varphi=14^0\,9'\,52''63$\\
mean daily sidereal motion&$=770''8398$\\
logarithm of the semimajor axis&$=0.4420283$\\
mean longitude at the first opposition&$=89^0\,20'\,20''65$
\end{tabular}\end{center}
\SourceLeaf{19}
\[
\begin{array}{rrr}
\gamma=-33^0&39'&35''63\\
\gamma'=+28&15&5.20\\
\gamma''=+27&20&9.32\\
\gamma'''=-4&52&52.65
\end{array}
\]

Comparison of the values of the heliocentric latitudes $\gamma$, $\gamma'$, $\gamma''$, $\gamma'''$ found in the three hypotheses yields
\[
\begin{array}{ll}
a=-0.0492,&b=+0.2087\\
a'=+0.0212,&b'=+0.0745\\
a''=+0.0317,&b''=+0.0042\\
a'''=-0.0112,&b'''=+0.0223
\end{array}
\]
We thus obtain four equations
\[
\begin{gathered}
+62''48-0.1744\mathrm d\AscendingNode+1.1957\mathrm di=0\\
+9.49+0.3578\mathrm d\AscendingNode-0.8172\mathrm di=0\\
-46.79-0.3292\mathrm d\AscendingNode-0.8685\mathrm di=0\\
-25.55-0.6915\mathrm d\AscendingNode+0.2034\mathrm di=0
\end{gathered}
\]
from which we deduce by the principle of least squares
\[
\begin{aligned}
\mathrm d\AscendingNode&=-54''41\\
\mathrm di&=-42.06
\end{aligned}
\]
\Needspace{6\baselineskip}
so that we obtain
\begin{center}\fontsize{9}{11}\selectfont\setlength{\tabcolsep}{3pt}\begin{tabular}{lr}
longitude of the ascending node for the beginning of 1803,&$172^0\,27'\,52''39$\\
orbital inclination $\ .\ .\ .\ .\ .\ .\ .\ .\ .\ .$&$34\quad36\quad49.44$
\end{tabular}\end{center}

The remaining elements can be obtained by comparing the results of the three hypotheses or, more accurately, by recalculating the longitudes in the orbit and repeating the operations explained in Article 5. The first method gives
\begin{center}\fontsize{9}{11}\selectfont\setlength{\tabcolsep}{3pt}\begin{tabular}{lr}
longitude of the perihelion 1803 $\ .\ .\ .\ .\ .\ .\ .$&$120^0\,58'\,3''86$\\
angle $\varphi\ .\ .\ .\ .\ .\ .\ .\ .\ .\ .\ .\ .\ .$&$14\quad9\quad36.25$\\
mean longitude at the first opposition $\ .\ .$&$89\quad20\quad32.08$\\
mean daily sidereal motion $\ .\ .\ .\ .\ .\ .\ .$&$770''7899$\\
logarithm of the semimajor axis. $\ .\ .\ .\ .\ ,\ .\ .$&$0.4420471$
\end{tabular}\end{center}
\SourceLeaf{20}
The other method gives
\begin{center}\fontsize{9}{11}\selectfont\setlength{\tabcolsep}{3pt}\begin{tabular}{lr}
longitude of the perihelion 1803 $\ .\ .\ .\ .\ .\ .\ .\ .$&$120^0\,58'\,4''81$\\
angle $\varphi\ .\ .\ .\ .\ .\ .\ .\ .\ .\ .\ .\ .\ .$&$14\quad9\quad36.63$\\
mean longitude at the first opposition $\ .\ .\ .\ .$&$89\quad20\quad31.81$\\
mean daily sidereal motion $\ .\ .\ .\ .\ .\ .\ .$&$770''7893$\\
logarithm of the semimajor axis $\ .\ .\ .\ .\ .\ .\ .\ .$&$0.4420473$
\end{tabular}\end{center}
These elements agree with those I gave above in Article 3.
\Article{9}
However great the perturbations Pallas undergoes from the other planets, experience shows that elliptic elements fitted to four oppositions match the planet’s motion fairly well throughout that interval. In fact they also differ little from its earlier and later motion, provided the time interval is not taken too large. For example, the second set of elements given in Article 3 differed from the observed heliocentric longitude by three minutes of arc at the opposition of 1803 and by five minutes at the opposition of 1809. This being so, in constructing an ephemeris for the planet’s future motion, it always seems most convenient to me to use purely elliptic elements derived from the four immediately preceding oppositions. The equations arising from perturbations will certainly become so numerous that calculating even one heliocentric position of the planet must turn out to be almost as laborious as calculating the elliptic elements by the instructions given above. One can thus always safely predict the geocentric position of the planet within an error of a few minutes of arc, which is precise enough to find it.
\Article{10}
Yet the dignity of science demands that we seek a more lasting agreement, which clearly cannot be obtained until perturbations are included in the calculation. In my judgment, such a lengthy and tedious calculation would have been premature as long as the available observations covered too short a time and the perturbations from the other planets scarcely made themselves felt. Now that elliptic motion no longer suffices to reconcile all the observed positions, the time seems to have come when we can think about a more accurate theory. I shall soon explain in more detail elsewhere how, in my view, to calculate the perturbations that Pallas undergoes, especially from Jupiter,
\SourceLeaf{21}
as conveniently and accurately as possible. The methods used for other planets can scarcely, if at all, be applied here because of the excessive eccentricity and inclination. Here I shall go on to discuss the elliptic elements that seem best suited as a basis for calculating the perturbations. My aim is to obtain elliptic elements that do not fit one or another of the oppositions exactly, but fit all those observed so far as closely as possible. I already gave a brief description of a method for carrying out this task in Article 187 of the Theoria Motus Corporum Coelestium. The things I treated there in general admit some shortcuts in the special case where the observed positions are oppositions. That work also leaves out some practical devices that I have long used to make the method of least squares easier to apply. I hope, therefore, that astronomers will welcome a somewhat fuller account of these calculations here. The whole task consists in determining the \textit{corrections} to be added to approximate elements that are assumed to differ little from all the observed positions. All the work thus depends on two steps. First the linear equations supplied by the individual observed positions must be formed. Then the most suitable values of the unknowns must be obtained from these equations.
\Article{11}
For the approximate elements, the symbols denote the following quantities.
\begin{description}\setlength{\itemsep}{0pt}\setlength{\parsep}{0pt}
\item[$L$] the mean longitude of the planet for an arbitrary epoch
\item[$t$] the number of days elapsed from the epoch to the instant of observation
\item[$\MeanMotionGlyph$] the mean daily sidereal motion in seconds of arc
\item[$\Pi$] the longitude of the perihelion
\item[$e=\sin\varphi$] the eccentricity
\item[$a$] the semimajor axis
\item[$r$] the radius vector
\item[$v$] the true anomaly
\item[$E$] the eccentric anomaly
\item[$\AscendingNode$] the longitude of the ascending node
\item[$i$] the orbital inclination
\item[$u$] the argument of latitude
\item[$\lambda$] the heliocentric longitude
\SourceLeaf{22}
\item[$\gamma$] the heliocentric latitude
\item[$\ApproxBeta$] the geocentric latitude
\item[$R$] the distance of the Earth from the Sun.
\end{description}
From the observation, however, I assume the following values.
\begin{description}\setlength{\itemsep}{0pt}\setlength{\parsep}{0pt}
\item[$\alpha$] the heliocentric longitude
\item[$\ObservedBeta$] the geocentric latitude.
\end{description}
Finally, by $\mathrm dL$, $\mathrm d\MeanMotionGlyph$, $\mathrm d\Pi$ and so on I denote the corrections to the quantities $L$, $\MeanMotionGlyph$, $\Pi$ and so on. Thus $\mathrm dL+t\mathrm d\MeanMotionGlyph$ is the correction to the mean longitude,
$\mathrm dL+t\mathrm d\MeanMotionGlyph-\mathrm d\Pi$ is the correction to the mean anomaly,
and hence by Articles 15 and 16 of the Theoria Motus Corp. Coel.
\[
\begin{aligned}
\mathrm dv&=\frac{aa\cos\varphi}{rr}(\mathrm dL+t\mathrm d\MeanMotionGlyph-\mathrm d\Pi)+\frac{aa}{rr}(2-e\cos E-ee)\sin E\mathrm d\varphi\\
\mathrm dr&=\frac ra\mathrm da+a\operatorname{tang}\varphi\sin v(\mathrm dL+t\mathrm d\MeanMotionGlyph-\mathrm d\Pi)-a\cos\varphi\cos v\mathrm d\varphi
\end{aligned}
\]
Furthermore, the correction to the argument of latitude is $\mathrm du=\mathrm dv+\mathrm d\Pi-\mathrm d\AscendingNode$,
and by Article 52 of the Theoria Motus C. C. the correction to the heliocentric longitude is given by
\[
\mathrm d\lambda=\mathrm d\AscendingNode-\operatorname{tang}\gamma\cos(\lambda-\AscendingNode)\mathrm di+\frac{\cos i}{\cos\gamma^2}\mathrm du
\]
From this we obtain
\[
\begin{aligned}
\mathrm d\lambda={}&\frac{aa\cos\varphi\cos i}{rr\cos\gamma^2}\mathrm dL\\
&+\frac{taa\cos\varphi\cos i}{rr\cos\gamma^2}\mathrm d\MeanMotionGlyph\\
&+\left(\frac{\cos i}{\cos\gamma^2}-\frac{aa\cos\varphi\cos i}{rr\cos\gamma^2}\right)\mathrm d\Pi\\
&+\frac{aa\cos i}{rr\cos\gamma^2}(2-e\cos E-ee)\sin E\mathrm d\varphi\\
&+\left(1-\frac{\cos i}{\cos\gamma^2}\right)\mathrm d\AscendingNode\\
&-\operatorname{tang}\gamma\cos(\lambda-\AscendingNode)\mathrm di
\end{aligned}
\]
Furthermore, since we have
\[
\begin{aligned}
a^{\frac32}\MeanMotionGlyph&=\mathrm{Const.}\\
r\sin(\ApproxBeta-\gamma)&=R\sin\ApproxBeta\\
\operatorname{tang}\gamma&=\operatorname{tang}i\sin(\alpha-\AscendingNode)
\end{aligned}
\]
differentiation gives
\SourceLeaf{23}
\[
\frac{\mathrm da}{a}=-\frac{2\mathrm d\MeanMotionGlyph}{3\MeanMotionGlyph}
\]
\[
\frac{\mathrm dr}{r}+\operatorname{cotang}(\ObservedBeta-\gamma).(\mathrm d\ApproxBeta-\mathrm d\gamma)=\operatorname{cotg}\ApproxBeta\mathrm d\ApproxBeta,\quad\text{or}
\]
\[
\mathrm d\ApproxBeta=\frac{\sin\ApproxBeta\cos(\ApproxBeta-\gamma)}{\sin\gamma}\mathrm d\gamma-\frac{\sin\ApproxBeta\sin(\ApproxBeta-\gamma)}{r\sin\gamma}\mathrm dr
\]
\[
\mathrm d\gamma=\frac{\sin2\gamma}{\sin2i}\mathrm di-\tfrac12\sin2\gamma\operatorname{cotang}(\alpha-\AscendingNode)\mathrm d\AscendingNode
\]
from which, using the expression for $\mathrm dr$ developed above, we obtain
\[
\begin{aligned}
\mathrm d\ApproxBeta={}&-\frac{a\sin\ApproxBeta\sin(\ApproxBeta-\gamma)\operatorname{tang}\varphi\sin v}{r\sin\gamma}\mathrm dL\\
&+\left\{\frac{2\sin\ApproxBeta\sin(\ApproxBeta-\gamma)}{3\MeanMotionGlyph\sin\gamma}-\frac{at\sin\ApproxBeta\sin(\ApproxBeta-\gamma)\operatorname{tang}\varphi\sin v}{r\sin\gamma}\right\}\mathrm d\MeanMotionGlyph\\
&+\frac{a\sin\ApproxBeta\sin(\ApproxBeta-\gamma)\operatorname{tang}\varphi\sin v}{r\sin\gamma}\mathrm d\Pi\\
&+\frac{a\sin\ApproxBeta\sin(\ApproxBeta-\gamma)\cos\varphi\cos v}{r\sin\gamma}\mathrm d\varphi\\
&+\frac{2\sin\ApproxBeta\cos(\ApproxBeta-\gamma)\cos\gamma}{\sin2i}\mathrm di\\
&-\sin\ApproxBeta\cos(\ApproxBeta-\gamma)\cos\gamma\operatorname{cotang}(\alpha-\AscendingNode)\mathrm d\AscendingNode
\end{aligned}
\]
The heliocentric longitude and geocentric latitude obtained from the corrected values of the elements are therefore $\lambda+\mathrm d\lambda$, $\ApproxBeta+\mathrm d\ApproxBeta$, and so each opposition supplies two equations
\[
\begin{aligned}
\alpha&=\lambda+\mathrm d\lambda\\
\ObservedBeta&=\ApproxBeta+\mathrm d\ApproxBeta
\end{aligned}
\]
\Article{12}
Applying these instructions to the six oppositions of Pallas given in Article 2,
and basing the calculation on the second set of elements outlined in Article 3, we obtain the following twelve equations.

\noindent From the \textit{first} opposition, where the longitude was found by calculation to be $=277^0\,36'\,20''07$,
\[
\text{geocentric latitude}=+46^0\,26'\,29''19:
\]
\[
\begin{aligned}
0={}&-183''93+0.79363\mathrm dL+143.66\mathrm d\MeanMotionGlyph+0.39493\mathrm d\Pi+0.95920\mathrm d\varphi\\
&-0.18856\mathrm d\AscendingNode+0.17387\mathrm di\\[4pt]
0={}&-6''81-0.02658\mathrm dL+46.71\mathrm d\MeanMotionGlyph+0.02658\mathrm d\Pi-0.20858\mathrm d\varphi\\
&+0.15946\mathrm d\AscendingNode+1.25782\mathrm di
\end{aligned}
\]
\SourceLeaf{24}
\Needspace{11\baselineskip}
\noindent From the \textit{second} opposition, where the calculated longitude is $=337^0\,0'\,36''04$,
\[
\text{geocentric latitude}=+15^0\,1'\,46''71:
\]
\[
\begin{aligned}
0={}&-0''06+0.58880\mathrm dL+358.12\mathrm d\MeanMotionGlyph+0.26208\mathrm d\Pi-0.85234\mathrm d\varphi\\
&+0.14912\mathrm d\AscendingNode+0.17775\mathrm di\\[4pt]
0={}&-3''09+0.01318\mathrm dL+28.39\mathrm d\MeanMotionGlyph-0.01318\mathrm d\Pi-0.07861\mathrm d\varphi\\
&+0.91704\mathrm d\AscendingNode+0.54365\mathrm di
\end{aligned}
\]
From the \textit{third} opposition, where the calculated longitude is $=67^0\,20'\,42''88$,
\[
\text{geocentric latitude}=-54^0\,31'\,3''88:
\]
\[
\begin{aligned}
0={}&-0''02+1.73436\mathrm dL+1846.17\mathrm d\MeanMotionGlyph-0.54603\mathrm d\Pi-2.05662\mathrm d\varphi\\
&-0.18833\mathrm d\AscendingNode-0.17445\mathrm di\\[4pt]
0={}&-8''98-0.12606\mathrm dL-227.42\mathrm d\MeanMotionGlyph+0.12606\mathrm d\Pi-0.38939\mathrm d\varphi\\
&+0.17176\mathrm d\AscendingNode-1.35441\mathrm di
\end{aligned}
\]
From the \textit{fourth} opposition, where the calculated longitude is \newline $=223^0\,37'\,25''39$,
\[
\text{geocentric latitude}=+42^0\,11'\,28''07:
\]
\[
\begin{aligned}
0={}&-2''31+0.99584\mathrm dL+1579.03\mathrm d\MeanMotionGlyph+0.06456\mathrm d\Pi+1.99545\mathrm d\varphi\\
&-0.06040\mathrm d\AscendingNode-0.33750\mathrm di\\[4pt]
0={}&+2''47-0.08089\mathrm dL-67.22\mathrm d\MeanMotionGlyph+0.08089\mathrm d\Pi-0.09970\mathrm d\varphi\\
&-0.46359\mathrm d\AscendingNode+1.22803\mathrm di
\end{aligned}
\]
From the \textit{fifth} opposition, where the calculated longitude is $=304^0\,2'\,59''71$,
\[
\text{geocentric latitude}=+37^0\,44'\,31''82:
\]
\[
\begin{aligned}
0={}&+0''01+0.65311\mathrm dL+1329.09\mathrm d\MeanMotionGlyph+0.38994\mathrm d\Pi-0.08439\mathrm d\varphi\\
&-0.04305\mathrm d\AscendingNode+0.34268\mathrm di\\[4pt]
0={}&+38''12-0.00218\mathrm dL+38.47\mathrm d\MeanMotionGlyph+0.00218\mathrm d\Pi-0.18710\mathrm d\varphi\\
&+0.47301\mathrm d\AscendingNode-1.14371\mathrm di
\end{aligned}
\]
From the \textit{sixth} opposition, where the calculated longitude is $=359^0\,34'\,46''67$,
\[
\text{geocentric latitude}=-7^0\,20'\,12''13:
\]
\[
\begin{aligned}
0={}&-317''73+0.69957\mathrm dL+1719.32\mathrm d\MeanMotionGlyph+0.12913\mathrm d\Pi-1.38787\mathrm d\varphi\\
&+0.17130\mathrm d\AscendingNode-0.08360\mathrm di\\[4pt]
0={}&+117''97-0.01315\mathrm dL-43.84\mathrm d\MeanMotionGlyph+0.01315\mathrm d\Pi+0.02929\mathrm d\varphi\\
&+1.02138\mathrm d\AscendingNode-0.27187\mathrm di
\end{aligned}
\]
Of these twelve equations, however, we shall reject the tenth altogether, since the observed geocentric latitude is too uncertain.
\SourceLeaf{25}
\Article{13}
Since we cannot determine the six unknowns $\mathrm dL$, $\mathrm d\MeanMotionGlyph$ and so on in such a way that all eleven equations are satisfied exactly, that is, that each of the functions of the unknowns on the right is simultaneously $=0$, we shall find the values for which the sum of the squares of these functions is as small as possible. It is easy to see that if the following linear functions of the unknowns $p$, $q$, $r$, $s$ and so on are given in general
\[
\begin{aligned}
n+ap+bq+cr+ds+\text{etc.}&=w\\
n'+a'p+b'q+c'r+d's+\text{etc.}&=w'\\
n''+a''p+b''q+c''r+d''s+\text{etc.}&=w''\\
n'''+a'''p+b'''q+c'''r+d'''s+\text{etc.}&=w'''
\end{aligned}
\]
etc., the equations of condition for $ww+w'w'+w''w''+w'''w'''+\text{etc.}=\Omega$ to be a minimum are these
\[
\begin{aligned}
aw+a'w'+a''w''+a'''w'''+\text{etc.}&=0\\
bw+b'w'+b''w''+b'''w'''+\text{etc.}&=0\\
cw+c'w'+c''w''+c'''w'''+\text{etc.}&=0\\
dw+d'w'+d''w''+d'''w'''+\text{etc.}&=0\\
\text{etc.\quad or}&
\end{aligned}
\]
using the following short notation
\[
\begin{array}{ll}
an+a'n'+a''n''+a'''n'''+\text{etc.}&\text{by }[an]\\
aa+a'a'+a''a''+a'''a'''+\text{etc.}&\text{by }[aa]\\
ab+a'b'+a''b''+a'''b'''+\text{etc.}&\text{by }[ab]\\
\text{etc.}&\\
bb+b'b'+b''b''+b'''b'''+\text{etc.}&\text{by }[bb]\\
bc+b'c'+b''c''+b'''c'''+\text{etc.}&\text{by }[bc]\\
\text{etc. etc.}&
\end{array}
\]
$p$, $q$, $r$, $s$ and so on must be determined by elimination from the equations
\[
\begin{aligned}{}
[an]+[aa]p+[ab]q+[ac]r+[ad]s+\text{etc.}&=0\\
[bn]+[ab]p+[bb]q+[bc]r+[bd]s+\text{etc.}&=0\\
[cn]+[ac]p+[bc]q+[cc]r+[cd]s+\text{etc.}&=0\\
[dn]+[ad]p+[bd]q+[cd]r+[dd]s+\text{etc.}&=0\\
\text{etc.}&
\end{aligned}
\]
\SourceLeaf{26}
Whenever the number of unknowns $p$, $q$, $r$, $s$ and so on is somewhat larger, however, elimination requires a very lengthy and tedious calculation, which can be considerably shortened as follows. In addition to the coefficients $[an]$, $[aa]$, $[ab]$ and so on, whose number is $=\tfrac12(ii+3i)$ if the number of unknowns is $=i$, I also assume that the following sum has been calculated, $nn+n'n'+n''n''+n'''n'''+\text{etc.}=[nn]$. It is then easy to see that
\[
\begin{aligned}
\Omega={}&[nn]+2[an]p+2[bn]q+2[cn]r+2[dn]s+\text{etc.}\\
&+[aa]pp+2[ab]pq+2[ac]pr+2[ad]ps+\text{etc.}\\
&\qquad+[bb]qq+2[bc]qr+2[bd]qs+\text{etc.}\\
&\qquad\qquad+[cc]rr+2[cd]rs+\text{etc.}\\
&\qquad\qquad\qquad+[dd]ss+\text{etc.}\\
&\qquad\qquad\qquad\qquad\text{etc.}
\end{aligned}
\]
If we therefore denote
\[
[an]+[aa]p+[ab]q+[ac]r+[ad]s+\text{etc.}\quad\text{by}\quad A
\]
it is clear that all the parts of $\dfrac{A^2}{[aa]}$ involving the factor $p$ are contained in $\Omega$, so that $\Omega-\dfrac{A^2}{[aa]}$ is a function independent of $p$. Hence, setting
\[
\begin{aligned}{}
[nn]-\frac{[an]^2}{[aa]}&=[nn,1]\\
[bn]-\frac{[an].[ab]}{[aa]}&=[bn,1]\\
[cn]-\frac{[an].[ac]}{[aa]}&=[cn,1]\\
[dn]-\frac{[an].[ad]}{[aa]}&=[dn,1]\quad\text{etc.}\\
[bb]-\frac{[ab]^2}{[aa]}&=[bb,1]\\
[bc]-\frac{[ab].[ac]}{[aa]}&=[bc,1]\\
[bd]-\frac{[ab].[ad]}{[aa]}&=[bd,1]\quad\text{etc. etc.: we get}
\end{aligned}
\]
\[
\begin{aligned}
\Omega-\frac{A^2}{[aa]}={}&[nn,1]+2[bn,1]q+2[cn,1]r+2[dn,1]s+\text{etc.}\\
&+[bb,1]qq+2[bc,1]qr+2[bd,1]qs+\text{etc.}\\
&\qquad+[cc,1]rr+2[cd,1]rs+\text{etc.}\\
&\qquad\qquad+[dd,1]ss+\text{etc.}\\
&\qquad\qquad\qquad\text{etc.}
\end{aligned}
\]
which is the function we shall denote by $\Omega'$.

\SourceLeaf{27}
Similarly, setting
\[
[bn,1]+[bb,1]q+[bc,1]r+[bd,1]s+\text{etc.}=B
\]
makes $\Omega'-\dfrac{B^2}{[bb,1]}$ a function independent of $q$, which we shall set $=\Omega''$. In the same way we set
\[
\begin{aligned}{}
[nn,1]-\frac{[bn,1]^2}{[bb,1]}&=[nn,2]\\
[cn,1]-\frac{[bn,1].[bc,1]}{[bb,1]}&=[cn,2]\\
[cc,1]-\frac{[bc,1]^2}{[bb,1]}&=[cc,2]\\
\text{etc. etc.\quad and}&
\end{aligned}
\]
\[
[cn,2]+[cc,2]r+[cd,2]s+\text{etc.}=C
\]
so that $\Omega''-\dfrac{C^2}{[cc,2]}$ is a function independent of $r$ as well. We proceed in the same way until we reach a term in the sequence $\Omega$, $\Omega'$, $\Omega''$ and so on that is independent of all the unknowns. It will be $[nn,\mu]$ if the number of unknowns $p$, $q$, $r$, $s$ and so on is denoted by $\mu$. We thus have
\[
\Omega=\frac{A^2}{[aa]}+\frac{B^2}{[bb,1]}+\frac{C^2}{[cc,2]}+\frac{D^2}{[dd,3]}+\text{etc.}+[nn,\mu]
\]
Now since $\Omega=ww+w'w'+w''w''\text{ etc.}$ by its nature cannot take a negative value, it is easily shown that the divisors $[aa]$, $[bb,1]$, $[cc,2]$, $[dd,3]$ and so on must necessarily be \textit{positive} (for brevity, however, I do not give this proof in detail here). It follows at once that the minimum value of $\Omega$ is obtained by setting $A=0$, $B=0$, $C=0$, $D=0$ and so on. From these $(\mu)$ equations, then, we must determine the unknowns $p$, $q$, $r$, $s$ and so on. This can very easily be done in reverse order, since the last equation evidently involves just one unknown, the next to last two, and so on. At the same time this method has the advantage that it gives the minimum value of the sum $\Omega$ directly, since it is evidently $=[nn,\mu]$.
\Article{14}
Let us now apply these instructions to our example, where $p$, $q$, $r$, $s$ and so on are $\mathrm dL$, $\mathrm d\MeanMotionGlyph$, $\mathrm d\Pi$, $\mathrm d\varphi$, $\mathrm d\AscendingNode$, $\mathrm di$. After carrying out the calculation carefully, I found the following numerical values.
\SourceLeaf{28}
{\fontsize{9}{11}\selectfont
\[
\begin{array}{lll}
\begin{gathered}{}
[nn]=148848\\
[an]=-371.09\\
[bn]=-580104\\
[cn]=-113.45\\
[dn]=+268.53\\
[en]=+94.26\\
[fn]=-31.81\\
[aa]=+5.91569\\
[ab]=+7203.91\\
[ac]=-0.09344\\
[ad]=-2.28516\\
[ae]=-0.34664\\
[af]=-0.18194\\
[bb]=+10834225\\
[bc]=-49.06\\
[bd]=-3229.77\\
[be]=-198.64\\
[bf]=-143.05\\
[cc]=+0.71917\\
[cd]=+1.13382\\
[ce]=+0.06400\\
[cf]=+0.26341\\
[dd]=+12.00340\\
[de]=-0.37137\\
[df]=-0{,}11762\\
[ee]=+2.28215\\
[ef]=-0.36136\\
[ff]=+5.62456\\[3pt]
\text{I then deduced}\\[3pt]
[nn,1]=+125569\\
[bn,1]=-138534
\end{gathered}&
\begin{gathered}{}
[cn,1]=-119.31\\
[dn,1]=-125.18\\
[en,1]=+72.52\\
[fn,1]=-43.22\\
[bb,1]=+2458225\\
[bc,1]=+62.13\\
[bd,1]=-510.58\\
[be,1]=+213.84\\
[bf,1]=+73.45\\
[cc,1]=+0.71769\\
[cd,1]=+1.09773\\
[ce,1]=-0.05852\\
[cf,1]=+0.26054\\
[dd,1]=+11.12064\\
[de,1]=-0.50528\\
[df,1]=-0.18790\\
[ee,1]=+2.26185\\
[ef,1]=-0.37202\\
[ff,1]=+5.61905\\[3pt]
\text{And hence, similarly}\\[3pt]
[nn,2]=+117763\\
[cn,2]=-115.81\\
[dn,2]=-153.95\\
[en,2]=+84.57\\
[fn,2]=-39.08\\
[cc,2]=+0.71612\\
[cd,2]=+1.11063\\
[ce,2]=-0.06392\\
[cf,2]=+0.25868\\
[dd,2]=+11.01463\\
[de,2]=-0.46088
\end{gathered}&
\begin{gathered}{}
[df,2]=-0.17265\\
[ee,2]=+2.24325\\
[ef,2]=-0.37841\\
[ff,2]=+5.61686\\[3pt]
\text{And from this}\\[3pt]
[nn,3]=+99034\\
[dn,3]=+25.66\\
[en,3]=+74.23\\
[fn,3]=+2.75\\
[dd,3]=+9.29213\\
[de,3]=-0.36175\\
[df,3]=-0.57384\\
[ee,3]=+2.23754\\
[ef,3]=-0.35532\\
[ff,3]=+5.52342\\[3pt]
\text{Hence, in the same way}\\[3pt]
[nn,4]=+98963\\
[en,4]=+75.23\\
[fn,4]=+4.33\\
[ee,4]=+2.22346\\
[ef,4]=-0.37766\\
[ff,4]=+5.48798\\[3pt]
\text{Hence}\\[3pt]
[nn,5]=+96418\\
[fn,5]=+17.11\\
[ff,5]=+5.42383\\[3pt]
\text{And hence, finally}\\[3pt]
[nn,6]=+96364
\end{gathered}
\end{array}
\]
}
We thus have the following six equations for determining the unknowns.
\[
\begin{aligned}
0={}&+17''11+5.42383\mathrm di\\
0={}&+75''23+2.22346\mathrm d\AscendingNode-0.37766\mathrm di\\
0={}&+25''66+9.29213\mathrm d\varphi-0.36175\mathrm d\AscendingNode-0.57384\mathrm di\\
0={}&-115''81+0.71612\mathrm d\Pi+1.11063\mathrm d\varphi-0.06392\mathrm d\AscendingNode+0.25868\mathrm di\\
0={}&-13854''+2458225\mathrm d\MeanMotionGlyph+62.13\mathrm d\Pi-510.58\mathrm d\varphi+213{,}84\mathrm d\AscendingNode\\
&+73.45\mathrm di\\[4pt]
0={}&-371''09+5.91569\mathrm dL+7203.91\mathrm d\MeanMotionGlyph-0.09344\mathrm d\Pi-2.28516\mathrm d\varphi\\
&-0.34664\mathrm d\AscendingNode-0.18194\mathrm di
\end{aligned}
\]
\SourceLeaf{29}
from which we obtain
\[
\begin{aligned}
\mathrm di&=-3''15\\
\mathrm d\AscendingNode&=-34''37\\
\mathrm d\varphi&=-4''29\\
\mathrm d\Pi&=+166''44\\
\mathrm d\MeanMotionGlyph&=+0''054335\\
\mathrm dL&=-3''06
\end{aligned}
\]

The corrected elliptic elements that fit all six oppositions as closely as possible are therefore the following.
\begin{center}\fontsize{9}{11}\selectfont\setlength{\tabcolsep}{3pt}\begin{tabular}{lr}
Epoch of mean longitude 1803, at the Göttingen meridian&$221^0\,34'\,53''64$\\
Mean tropical daily motion $\ .\ .\ .\ .\ .\ .\ .\ .\ .\ .\ .\ .$&$770''5010$\\
Longitude of perihelion 1803 $\ .\ .\ .\ .\ .\ .\ .\ .\ .\ .\ .\ .\ .$&$121^0\,8'\,8''54$\\
Longitude of ascending node 1803 $\ .\ .\ .\ .\ .\ .\ .\ .\ .\ .$&$172\quad28\quad12.43$\\
Inclination of orbit $\ .\ .\ .\ .\ .\ .\ .\ .\ .\ .\ .\ .\ .\ .\ .$&$34\quad37\quad28.35$\\
Eccentricity $=\sin14^0\,9'\,59''79\ .\ .\ .\ .\ .\ .\ .\ .\ .\ .\ .$&$0.2447424$\\
Logarithm of semimajor axis $\ .\ .\ .\ .\ .\ .\ .\ .\ .\ .\ .\ .$&$0.4422071$
\end{tabular}\end{center}
\Article{15}
Substituting the values just found for the corrections $\mathrm dL$, $\mathrm d\MeanMotionGlyph$ and so on into the twelve equations of Article 12, we obtain the following differences between the observed and computed values of the heliocentric longitudes and geocentric latitudes.
\begin{center}\begin{tabular}{c|r|r}
At the opposition&\multicolumn{2}{c}{Difference}\\
of the year&\multicolumn{1}{c|}{in longitude}&\multicolumn{1}{c}{in latitude}\\\hline
1803&$-111''00$&$-8''31$\\
1804&$+59.18$&$-36.67$\\
1805&$+19.92$&$+0.07$\\
1807&$+85.77$&$+25.01$\\
1808&$+135.88$&$+28.72$\\
1809&$-216.54$&$+83.01$
\end{tabular}\end{center}
